\usepackage{graphicx}
\usepackage{newunicodechar}
\usepackage{pdflscape}
\usepackage[most]{tcolorbox}
\usepackage{xurl}
\setkeys{Gin}{keepaspectratio}
\newfontfamily\documentsymbolfont{Noto Sans Symbols}
\newunicodechar{−}{\ensuremath{-}}
\newunicodechar{≤}{\ensuremath{\leq}}
\newunicodechar{≥}{\ensuremath{\geq}}
\newunicodechar{♀}{{\documentsymbolfont\char"2640}}
\newunicodechar{♂}{{\documentsymbolfont\char"2642}}
\setlength{\emergencystretch}{3em}
\providecommand{\documentationnotetitle}{Note}
\providecommand{\documentationtiptitle}{Tip}
\providecommand{\documentationimportanttitle}{Important}
\providecommand{\documentationwarningtitle}{Warning}
\providecommand{\documentationcautiontitle}{Caution}
\providecommand{\documentationplannedtitle}{Planned}
\newtcolorbox{documentationnote}{
  breakable,
  colback=documentationnotebackground,
  colframe=documentationnoteborder,
  fonttitle=\bfseries,
  title=\documentationnotetitle
}
\newtcolorbox{documentationtip}{
  breakable,
  colback=documentationtipbackground,
  colframe=documentationtipborder,
  fonttitle=\bfseries,
  title=\documentationtiptitle
}
\newtcolorbox{documentationimportant}{
  breakable,
  colback=documentationimportantbackground,
  colframe=documentationimportantborder,
  fonttitle=\bfseries,
  title=\documentationimportanttitle
}
\newtcolorbox{documentationwarning}{
  breakable,
  colback=documentationwarningbackground,
  colframe=documentationwarningborder,
  fonttitle=\bfseries,
  title=\documentationwarningtitle
}
\newtcolorbox{documentationcaution}{
  breakable,
  colback=documentationcautionbackground,
  colframe=documentationcautionborder,
  fonttitle=\bfseries,
  title=\documentationcautiontitle
}
\newtcolorbox{documentationplanned}{
  breakable,
  colback=documentationplannedbackground,
  colframe=documentationplannedborder,
  fonttitle=\bfseries,
  title=\documentationplannedtitle
}
% Code is typeset by Pandoc's fancyvrb path (Shaded/Highlighting) and by the
% plain verbatim environment below, not by listings: under XeLaTeX listings
% cannot process characters above U+00FF, so Czech letters, arrows and box
% drawing in code were reordered, set in the wrong font or broke the build.
\usepackage{fvextra}
\definecolor{documentationcode}{HTML}{435489}
\definecolor{documentationcodebackground}{HTML}{F7F7F7}
\definecolor{documentationcoderule}{HTML}{B3B2B3}
\tcbset{documentationcodebox/.style={breakable,enhanced,arc=0pt,boxrule=0.4pt,
  colback=documentationcodebackground,colframe=documentationcoderule,
  left=4pt,right=4pt,top=2pt,bottom=2pt,before skip=0.6em,after skip=0.8em}}
\renewenvironment{verbatim}{%
  \VerbatimEnvironment
  \begin{tcolorbox}[documentationcodebox]%
  \begin{Verbatim}[breaklines,breakanywhere,fontsize=\small]%
}{%
  \end{Verbatim}%
  \end{tcolorbox}%
}
\AtBeginDocument{%
  \ifcsname Shaded\endcsname
    \renewenvironment{Shaded}{\begin{tcolorbox}[documentationcodebox]}{\end{tcolorbox}}%
  \fi
  \let\documentationoriginaltexttt\texttt
  \renewcommand{\texttt}[1]{{\color{documentationcode}\documentationoriginaltexttt{#1}}}%
}
% Arrows, box drawing, geometric shapes and check marks are missing from the
% theme fonts. They come from DejaVu Sans in text and DejaVu Sans Mono in code.
\newfontfamily\documentfallbackfont{DejaVu Sans}
\newfontfamily\documentfallbackmonofont{DejaVu Sans Mono}
\makeatletter
\newcount\documentation@fallbackchar
\def\documentation@fallbackglyph#1{%
  \edef\documentation@currentfamily{\f@family}%
  \edef\documentation@monofamily{\ttdefault}%
  \ifx\documentation@currentfamily\documentation@monofamily
    {\documentfallbackmonofont\char#1\relax}%
  \else
    {\documentfallbackfont\char#1\relax}%
  \fi}
\def\documentation@fallbackrange#1#2{%
  \documentation@fallbackchar=#1\relax
  \loop
    \begingroup
      \lccode`\~=\documentation@fallbackchar
      \lowercase{\endgroup
        \edef~{\noexpand\documentation@fallbackglyph{\the\documentation@fallbackchar}}}%
    \catcode\documentation@fallbackchar=13\relax
  \ifnum\documentation@fallbackchar<#2\relax
    \advance\documentation@fallbackchar by 1\relax
  \repeat}
\documentation@fallbackrange{"2190}{"21FF}
\documentation@fallbackrange{"2500}{"257F}
\documentation@fallbackrange{"25A0}{"25FF}
\documentation@fallbackrange{"2713}{"2718}
\makeatother
\RedeclareSectionCommand[tocnumwidth=5.75em]{subsection}
\RedeclareSectionCommand[tocnumwidth=7em]{subsubsection}
